\documentclass[11pt,a4paper]{article}
\usepackage[margin=23mm]{geometry}
\usepackage[T1]{fontenc}
\usepackage{amsmath,amssymb,amsthm,booktabs,tabularx}
\usepackage[colorlinks=true,linkcolor=blue,citecolor=blue,urlcolor=blue]{hyperref}
\hypersetup{pdftitle={Boundary formulas for Hardin's array recurrences and Barker's generating functions via Biggs--Meredith enumeration},pdfauthor={Haobo Ma; Wenlin Zhang}}
\newtheorem{theorem}{Theorem}
\newtheorem{corollary}[theorem]{Corollary}
\newcommand{\ones}{\boldsymbol{1}}
\newcommand{\counts}{C}
\setlength{\emergencystretch}{1em}
\title{Boundary formulas for Hardin's array recurrences and Barker's generating functions via Biggs--Meredith enumeration}
\author{\begin{tabular}{c}Haobo Ma\\[-1mm]\small ChronoAI Pte Ltd, Singapore\end{tabular}
\and \begin{tabular}{c}Wenlin Zhang\\[-1mm]\small National University of Singapore\\[-1mm]\small ChronoAI Pte Ltd, Singapore\end{tabular}}
\date{19 November 2026}

\begin{document}
\maketitle
\begin{abstract}
We give an alternative proof of Hardin's empirical recurrence and Barker's
generating function in OEIS A233156, revision 8 (9 October 2018), FORMULA
lines 1 and 2, by reducing them to Biggs and Meredith's double-strip
enumeration (1976, Section~3, pp.~10--11).
For labeled three-row arrays over \(\{0,1,2\}\) whose horizontal and
antidiagonal neighbors never sum to two, the count is exactly 36 times
the number of proper three-colorings of a three-row rectangular grid
with two fewer columns. This gives \(a_n=5a_{n-1}-2a_{n-2}\) for
\(n\ge5\) and the complete generating function, including its exceptional
initial terms. We also give a five-state proof, establish sharp onset and
minimal eventual order two, and prove every displayed row and column
recurrence in the parent table A233155, including Barker's four-row
generating function. The larger strips use specified finite integer
identities.
\end{abstract}

\section{The counting problem}
For positive integers \(h,n\), let \(\counts_h(n)\) be the number of arrays
\(x=(x_{r,j})\in\{0,1,2\}^{h\times n}\) such that
\begin{align}
 x_{r,j}+x_{r,j+1}&\ne2 &&(1\le r\le h,\ 1\le j<n),\label{eq:horizontal}\\
 x_{r+1,j}+x_{r,j+1}&\ne2 &&(1\le r<h,\ 1\le j<n).\label{eq:diagonal}
\end{align}
There is no restriction between vertically adjacent entries. Every labeled
array counts separately. This is the table A233155 of Hardin~\cite{parent};
its third row is A233156~\cite{source}. Write \(a_n=\counts_3(n)\).
The first FORMULA line of A233156 calls its recurrence ``Empirical'';
the second, signed by Colin Barker, gives an ``Empirical g.f.''.
The theorem below establishes both formulas as consequences of classical
grid enumeration, with the boundary factors made explicit.

The general transfer-matrix method produces rational counting series;
see~\cite[Section V.6]{ac}.
Biggs and Meredith~\cite[Section~3, Theorem~2 and p.~11]{bm} already give
the relevant three-row grid recurrence. Ro\v{c}ek, Shrock and Tsai~\cite[
Section~II.B]{rst} give its rational generating function; Salas and
Sokal~\cite[Section~5.2]{ss} give a transfer matrix and explicitly credit
Biggs and Meredith. A later combinatorial proof is supplied by Ginepro
and Hull~\cite[Section~3 and Theorem~4]{gh}, through their bijection with
Miura-ori assignments; Spahn~\cite[Section~6.1]{spahn} gives a further
finite-state derivation. The connection here is an exact boundary
reduction, not just a match of recurrence denominators. The five-state
argument below also proves the result directly from the array definition.

\begin{theorem}\label{thm:main}
The initial values are
\[
 a_1=27,\qquad a_2=96,\qquad a_3=432,\qquad a_4=1944.
\]
For every \(n\ge5\),
\begin{equation}\label{eq:recurrence}
 a_n=5a_{n-1}-2a_{n-2}.
\end{equation}
Consequently the ordinary generating function, starting at \(n=1\), is
\begin{equation}\label{eq:gf}
 \sum_{n\ge1}a_nz^n
 =\frac{3z(9-13z+2z^2-8z^3)}{1-5z+2z^2}.
\end{equation}
Equation~\eqref{eq:recurrence} fails at \(n=4\), and no constant-coefficient
order-one homogeneous recurrence holds eventually.
\end{theorem}

\section{The coloring graph and its rectangular core}\label{sec:boundary}
Let \(\tau(t)=2-t\), an involution of \(\{0,1,2\}\), and set
\(y_{r,j}=\tau^{j-1}(x_{r,j})\).
Every edge in~\eqref{eq:horizontal}--\eqref{eq:diagonal} joins consecutive
columns. For its two labels \(u,v\),
\[
 \tau^{j-1}(u)\ne\tau^j(v)
 \quad\Longleftrightarrow\quad u\ne\tau(v)
 \quad\Longleftrightarrow\quad u+v\ne2.
\]
Thus this relabeling is a bijection from the original arrays to proper
three-colorings of the graph with exactly those edges. It is invertible
column by column and introduces no normalization by color permutations.
Let \(G_n\) denote this coloring graph when \(h=3\), and let \(P_m\)
denote the path on \(m\) vertices. Write
\[
 b_m=P(P_3\mathbin{\square}P_m,3)\qquad(m\ge1),
\]
where \(\square\) is the Cartesian product of graphs and \(P(H,q)\)
is the chromatic polynomial of a finite graph \(H\), counting proper
colorings by \(q\) labeled colors for positive integers \(q\).

\begin{theorem}\label{thm:boundary}
For every \(n\ge3\),
\begin{equation}\label{eq:boundary}
 P(G_n,q)=(q-1)^2(q^2-3q+3)^2
 P(P_3\mathbin{\square}P_{n-2},q).
\end{equation}
In particular, \(a_n=36b_{n-2}\).
\end{theorem}
\begin{proof}
Shear the coordinates by \((r,j)\mapsto(r,t)=(r,r+j)\).
Horizontal edges become \((r,t)\sim(r,t+1)\), and antidiagonal edges
become \((r,t)\sim(r+1,t)\). The vertices
\[
 1\le r\le3,\qquad4\le t\le n+1
\]
therefore induce \(P_3\square P_{n-2}\).
Exactly three vertices remain on each side:
\[
 L=\{(1,2),(1,3),(2,3)\},\qquad
 R=\{(2,n+2),(3,n+2),(3,n+3)\}.
\]
On the left, \((1,3),(2,3)\) form, together with the core edge
\((1,4)\sim(2,4)\), an attached square; \((1,2)\) is pendant at
\((1,3)\). On the right there is the same configuration along the core
edge \((2,n+1)\sim(3,n+1)\), with pendant vertex \((3,n+3)\).
There are no other edges incident to \(L\cup R\), and the two sets are
disjoint even when \(n=3\).

Fix a proper coloring of the core. On an attached square the shared
edge has two distinct colors. Its two new vertices have \((q-1)^2\)
color choices before their mutual edge is imposed; exactly \(q-2\)
of these give them equal colors. Thus there are
\(q^2-3q+3\) square extensions, each with \(q-1\) choices for the
pendant vertex. The choices at the two ends are independent conditional
on the core coloring. This proves~\eqref{eq:boundary} for all integers
\(q\ge2\), and hence as a polynomial identity. At \(q=3\) each end
has six extensions, proving the factor 36.
\end{proof}

\subsection*{The classical count and the full numerator}
The double-strip matrix on p.~11 of Biggs and Meredith~\cite{bm},
specialized to three colors, is
\[
 T=\begin{pmatrix}4&1\\2&1\end{pmatrix},\qquad
 T^2-5T+2I=0.
\]
Their initial graph is a three-vertex path. The two components count
the strip and the quotient identifying the top and bottom boundary
vertices. Initially these counts are \(12\) and \(6\); after one
step the first is \(4(12)+6=54\).
Theorem~2 of that paper therefore gives
\[
 b_1=12,\quad b_2=54,\qquad
 b_m=5b_{m-1}-2b_{m-2}\quad(m\ge3).
\]
Equivalently,
\begin{equation}\label{eq:grid-gf}
 B(z):=\sum_{m\ge1}b_mz^m=\frac{12z-6z^2}{1-5z+2z^2}.
\end{equation}
Ro\v{c}ek--Shrock--Tsai~\cite[eqs.~(2.35)--(2.36),
(2.40)--(2.41)]{rst} give
\(\sum_{j\ge0}b_{j+2}z^j=(54-24z)/(1-5z+2z^2)\),
which is the same formula with a different starting width.
Ginepro--Hull's normalized three-row series~\cite[p.~6]{gh} is
\((1-z)/(1-5z+2z^2)\); their Theorem~4 identifies its positive-width
coefficients with \(b_m/3\), since one vertex is precolored.
Spahn's series~\cite[pp.~44--45]{spahn} is \(1+B(z)\).

Finally \(a_1=27\) and \(a_2=96\): one column is unrestricted, while
the two-column graph is a path on six vertices, with \(3\cdot2^5\)
colorings. Thus Theorem~\ref{thm:boundary} yields
\[
 \sum_{n\ge1}a_nz^n=27z+96z^2+36z^2B(z)
 =\frac{27z-39z^2+6z^3-24z^4}{1-5z+2z^2}.
\]
It also gives the recurrence for \(n\ge5\). The nonzero coefficient
of \(z^4\) proves that the cutoff is sharp. These are exact consequences
of the older grid count with the explicit boundary decomposition above.

\begin{corollary}\label{cor:oeis}
With the offsets in OEIS A020698 and A052913, for every \(n\ge3\),
\[
 a_n=216\,\mathrm{A020698}(n-3)
     =108\,\mathrm{A052913}(n-2).
\]
\end{corollary}
\begin{proof}
The defining recurrence and initial values \((2,9)\) of A020698
\cite{related-one} give \(b_m=6\,\mathrm{A020698}(m-1)\) for \(m\ge1\).
The defining recurrence and initial values \((1,4)\) of A052913
\cite{related-two} give \(b_m=3\,\mathrm{A052913}(m)\) for \(m\ge1\).
Both identities follow by their two initial values and induction.
Theorem~\ref{thm:boundary} proves the asserted equalities.
\end{proof}
These identities are not posed in the cited array entries. The first
also agrees with Andrew Woods's 7 September 2011 grid-coloring
interpretation in A020698, where all six permutations of the three
colors are identified.

\section{A direct five-state proof}

A column \((u_1,u_2,u_3)\) has one of the five equality patterns
\[
 AAA,\quad AAB,\quad ABA,\quad ABB,\quad ABC.
\]
These denote, respectively, three equal entries; only the first two equal;
only the first and third equal; only the last two equal; or three distinct
entries. The letters stand for distinct colors when different.
Their multiplicities among the 27 labeled columns form the row vector
\[
 \alpha=(3,6,6,6,6).
\]
A column \(u\) may be followed by \(v\) exactly when
\begin{equation}\label{eq:transition}
 v_1\notin\{u_1,u_2\},\qquad
 v_2\notin\{u_2,u_3\},\qquad v_3\ne u_3.
\end{equation}
Simultaneously permuting all three colors preserves these conditions and
the equality pattern of \(v\). Hence the number of successors in each
pattern depends only on the pattern of \(u\).

For completeness, the successor columns for representatives are
\[
\begin{array}{c|l}
000 & 111,112,121,122,211,212,221,222\\
001 & 120,122,220,222\\
010 & 221,222\\
011 & 200,202,220,222\\
012 & 200,201
\end{array}
\]
by~\eqref{eq:transition}. Grouping each row by its successor's pattern gives
\begin{equation}\label{eq:matrix}
 Q=\begin{pmatrix}
 2&2&2&2&0\\
 1&1&0&1&1\\
 1&1&0&0&0\\
 1&1&1&1&0\\
 0&0&0&1&1
 \end{pmatrix}.
\end{equation}
Let \(\ones=(1,1,1,1,1)^{\mathsf T}\). The entry of
\(Q^k\ones\) for a given pattern counts the ways to append \(k\) columns
to any fixed labeled column of that pattern. Indeed this holds for
\(k=0\); choosing the next column and then its continuation proves the
induction step. Therefore
\begin{equation}\label{eq:count}
 a_n=\alpha Q^{n-1}\ones\qquad(n\ge1).
\end{equation}
The initial multiplicities \(\alpha\), not the number of patterns, ensure
that all labeled arrays are counted.

Direct multiplication of~\eqref{eq:matrix} gives
\[
\begin{array}{c|rrrrr}
k & \multicolumn{5}{c}{(Q^k\ones)^{\mathsf T}}\\\hline
0&1&1&1&1&1\\
1&8&4&2&4&2\\
2&36&18&12&18&6\\
3&168&78&54&84&24\\
4&768&354&246&384&108
\end{array}
\]
so, entry by entry,
\begin{equation}\label{eq:identity}
 Q^4\ones-5Q^3\ones+2Q^2\ones=0.
\end{equation}
Multiplying by \(\alpha Q^{n-5}\) proves~\eqref{eq:recurrence} for
every \(n\ge5\). Multiplying the first four rows of the table by
\(\alpha\) gives the stated initial values. Notice that this is an identity
on continuation vectors; it need not be an identity on the whole matrix.

For the generating function, take \(F(z)=\sum_{n\ge1}a_nz^n\) as a formal
power series. The coefficients of \((1-5z+2z^2)F(z)\) in degrees at least
five vanish by~\eqref{eq:recurrence}; those in degrees one through four are
\[
27,\quad 96-5(27)=-39,\quad
432-5(96)+2(27)=6,\quad
1944-5(432)+2(96)=-24.
\]
This is the numerator in~\eqref{eq:gf}. Its last coefficient also shows
that the recurrence does not extend to \(n=4\).

Finally, \(a_5=8856\), and put
\(D_n=a_na_{n+2}-a_{n+1}^2\) for \(n\ge3\).
The recurrence gives
\[
 D_3=432(8856)-1944^2=46656,\qquad D_{n+1}=2D_n\quad(n\ge3).
\]
For the second equality, replace \(a_{n+3}\) by
\(5a_{n+2}-2a_{n+1}\) in \(D_{n+1}\), and then replace
\(5a_{n+1}-a_{n+2}\) by \(2a_n\).
Thus \(D_n\ne0\) for every \(n\ge3\), whereas an eventual order-one
homogeneous recurrence would make all sufficiently late \(D_n\) zero.
This completes the proof of Theorem~\ref{thm:main}.

\section{Consequences and related entries}
\begin{corollary}\label{cor:closed}
Let \(\lambda_\pm=(5\pm\sqrt{17})/2\). For \(n\ge3\),
\[
 a_n=216\left[\left(1+\frac4{\sqrt{17}}\right)\lambda_+^{n-3}
             +\left(1-\frac4{\sqrt{17}}\right)\lambda_-^{n-3}\right].
\]
In particular, \(a_{n+1}/a_n\to\lambda_+\).
\end{corollary}
\begin{proof}
Both powers satisfy~\eqref{eq:recurrence}. The displayed expression is
\(432\) at \(n=3\) and \(1944\) at \(n=4\), so induction proves the
formula. Since \(0<\lambda_-<\lambda_+\) and both coefficients are
positive, division by \(\lambda_+^{n-3}\) gives the ratio limit.
\end{proof}

\begin{corollary}\label{cor:thin}
The parent table A233155 satisfies
\begin{align*}
 \counts_1(n)&=3\cdot2^{n-1} &&(n\ge1),\\
 \counts_2(1)&=9,\qquad
 \counts_2(n)=24\cdot3^{n-2} &&(n\ge2),\\
 \counts_h(1)&=3^h,\quad
 \counts_h(2)=6\cdot4^{h-1} &&(h\ge1),\\
 \counts_h(3)&=12\cdot6^{h-1},\quad
 \counts_h(4)=24\cdot9^{h-1} &&(h\ge1).
\end{align*}
\end{corollary}
\begin{proof}
Use the same alternating-column bijection for arbitrary height.
For one row, a first color has three choices and each successor has two.
For two rows the column patterns are \(AA,AB\), with multiplicities
\(\beta=(3,6)\) and transfer matrix
\[
 R=\begin{pmatrix}2&2\\1&1\end{pmatrix}.
\]
Indeed a column \(00\) has successors \(11,12,21,22\), and a column
\(01\) has successors \(20,22\). With \(\ones_2=(1,1)^{\mathsf T}\),
\(\counts_2(n)=\beta R^{n-1}\ones_2\),
\(\beta R\ones_2=24\), and \(R^2\ones_2=3R\ones_2\).
These identities give the two-row formula for all \(n\ge2\).

A single column has no edges, giving \(3^h\).
For two columns, the graph is the path
\[
 (1,1),(1,2),(2,1),(2,2),\ldots,(h,1),(h,2),
\]
with \(2h\) vertices and exactly the specified edges. Its colorings number
\(3\cdot2^{2h-1}=6\cdot4^{h-1}\).
For three columns, the first row is a properly colored path and has
\(3\cdot2^2=12\) choices. Suppose the preceding row is
\((u_1,u_2,u_3)\), where \(u_2\ne u_3\).
The next row \((v_1,v_2,v_3)\) must satisfy
\(v_1\ne u_2\), \(v_2\ne u_3\), \(v_1\ne v_2\), and \(v_2\ne v_3\).
There are two possibilities for \(v_2\): if \(v_2=u_2\), there are
two choices for \(v_1\); otherwise there is one. In each case \(v_3\)
has two choices. Thus every row has exactly \((2+1)2=6\) successors,
giving \(12\cdot6^{h-1}\).

For four columns, each proper row has one of the patterns
\(ABAB,ABAC,ABCA,ABCB\), each with six labeled realizations.
A fixed row \(u\) permits a proper row \(v\) below it exactly when
\(v_j\ne u_{j+1}\) for \(1\le j\le3\).
Counting the 24 possible \(v\)'s by pattern for representatives
\(0101,0102,0120,0121\) gives
\[
 S=\begin{pmatrix}3&3&2&2\\2&2&2&2\\2&2&2&2\\3&3&2&2\end{pmatrix}.
\]
Thus \(\counts_h(4)=(6,6,6,6)S^{h-1}\ones_4\), where
\(\ones_4=(1,1,1,1)^{\mathsf T}\).
Here \(S\ones_4=(10,8,8,10)^{\mathsf T}\) and
\(S^2\ones_4=9S\ones_4\).
The first two counts are \(24\) and \(216\), so the formula
\(\counts_h(4)=24\cdot9^{h-1}\) follows for every \(h\ge1\).
\end{proof}

\section{The remaining parent-table recurrences}\label{sec:siblings}
The FORMULA block of A233155, revision 6 (3 September 2022), lists
empirical recurrences for columns 1--7 and rows 1--7. The words
``[order 7]'', ``[order 9]'' and ``[order 27]'' in its last three rows
are made explicit in A233158, A233159 and A233160, respectively
\cite{higher}. The same representation establishes all these recurrences
and the two empirical FORMULA lines in A233157, revision 8
(9 October 2018), including Barker's generating function~\cite{four}.
The four-row square-strip denominator \(1-9z+15z^2-6z^3\)
already occurs in Ro\v{c}ek--Shrock--Tsai~\cite[eqs.~(3.4)--(3.6)]{rst}.
Their diagonal-strip denominators, at three colors, are
\(1-15z+18z^2\) and \(1-25z+90z^2-81z^3\)
\cite[Appendix~II.A, eqs.~(8.3)--(8.4), (8.8)--(8.10)]{rst}; these
also occur in columns five and six below. Different boundary graphs
can change the numerator and onset, so these denominator coincidences
alone do not identify the full counts. The following finite-set
construction supplies the actual array boundaries and the stated
initial ranges for every displayed sibling formula.

For a tuple \(u\) of colors, let \(\pi(u)\) replace the first color
encountered by 0, the next new color by 1, and the third by 2.
Let \(\mathcal P_h=\{\pi(u):u\in\{0,1,2\}^h\}\), in lexicographic
order. These are all equality patterns of unrestricted columns.
For \(u,v\in\mathcal P_h\), define
\[
 (Q_h)_{u,v}=\#\{w\in\{0,1,2\}^h:\pi(w)=v,\quad
 w_i\ne u_i\ (1\le i\le h),\quad
 w_i\ne u_{i+1}\ (1\le i<h)\}.
\]
Let \(\alpha_h(u)=3\) if \(u\) is constant and \(\alpha_h(u)=6\)
otherwise. Color permutation symmetry and the continuation induction
used in~\eqref{eq:count} give
\begin{equation}\label{eq:generic-row}
 \counts_h(n)=\alpha_h Q_h^{n-1}\ones_h\quad(n\ge1),
\end{equation}
where \(\ones_h\) is the all-one column indexed by \(\mathcal P_h\).
In particular \(Q_3=Q\); the quotient counts labeled arrays, not
color-permutation classes of whole arrays.

For a fixed width \(k\ge2\), let \(\mathcal R_k\) be the lexicographically
ordered patterns \(\pi(u)\) of proper rows, that is,
\(u_i\ne u_{i+1}\) for \(1\le i<k\). Set
\[
 (S_k)_{u,v}=\#\{w\in\{0,1,2\}^k:\pi(w)=v,\quad
 w_i\ne w_{i+1},\ w_i\ne u_{i+1}\ (1\le i<k)\}.
\]
Every proper row has six labeled realizations, so
\begin{equation}\label{eq:generic-column}
 \counts_h(k)=6\boldsymbol{1}^{\mathsf T}S_k^{h-1}\boldsymbol{1}
 \quad(h\ge1),
\end{equation}
with all-one vectors indexed by \(\mathcal R_k\).
This follows by the same induction, now appending rows. All matrices
in these definitions are specified by finite sets, without fitted entries.

\begin{corollary}\label{cor:all}
For the coefficient list \(c=(c_1,\ldots,c_d)\) in
Table~\ref{tab:rows}, the row sequence \(b_n=\counts_h(n)\) satisfies
\(b_n=\sum_{j=1}^d c_jb_{n-j}\) for every \(n\ge N\) in that table.
For \(k=5,6,7\), the column sequence \(b_h=\counts_h(k)\) satisfies,
respectively,
\begin{align*}
 b_h&=15b_{h-1}-18b_{h-2} &&(h\ge3),\\
 b_h&=25b_{h-1}-90b_{h-2}+81b_{h-3} &&(h\ge4),\\
 b_h&=42b_{h-1}-351b_{h-2}+972b_{h-3}-810b_{h-4} &&(h\ge5).
\end{align*}
For four rows the ordinary generating function is
\[
 \sum_{n\ge1}\counts_4(n)z^n=
 \frac{3z(27-115z+117z^2-186z^3+96z^4-360z^5+216z^6)}
 {1-9z+15z^2-6z^3}.
\]
Its recurrence begins sharply at \(n=8\).
\end{corollary}

\begin{table}[ht]
\centering\small
\begin{tabularx}{\textwidth}{@{}rrrX@{}}\toprule
Height \(h\)&\(N\)&\(d\)&\((c_1,\ldots,c_d)\)\\\midrule
3&5&2&\((5,-2)\)\\
4&8&3&\((9,-15,6)\)\\
5&12&7&\((16,-64,76,17,-84,48,-8)\)\\
6&16&9&\((30,-283,1038,-573,-6705,21056,-27636,17216,-4128)\)\\
7&34&27&\(\begin{aligned}[t](&55,-1066,8965,-19778,-225541,2020719,\\
 &-6851286,7381301,20655503,-79401996,80882723,\\
 &65337447,-222811374,135953338,120111164,\\
 &-197698384,43440564,76879396,-52839536,\\
 &-1810608,13170096,-4166800,-433792,532928,\\
 &-126464,13056,-512)\end{aligned}\)\\\bottomrule
\end{tabularx}
\caption{The empirical row recurrences of the parent table, with
the coefficient lists completed from their individual entries.}
\label{tab:rows}
\end{table}

\begin{proof}
Write \(p(X)=X^d-\sum_{j=1}^dc_jX^{d-j}\).
The finite identities required for the rows are
\begin{equation}\label{eq:row-cert}
 Q_h^s p(Q_h)\ones_h=0,
 \qquad (h,s)=(3,2),(4,4),(5,4),(6,6),(7,6).
\end{equation}
The dimensions of these matrices are respectively \(5,14,41,122,365\).
For columns 5, 6 and 7, use the corresponding polynomial from the
displayed column recurrences. The finite identities are
\begin{equation}\label{eq:column-cert}
 p(S_5)\boldsymbol{1}=0,\qquad
 S_6p(S_6)\boldsymbol{1}=0,\qquad
 p(S_7)\boldsymbol{1}=0;
\end{equation}
the dimensions are 8, 16 and 32. The complete matrices and their
successive vectors \(v_0=\boldsymbol{1}\), \(v_{j+1}=Mv_j\), through
\(v_{s+d}\), are given in \texttt{certificates.json}.
Each identity is the exact integer calculation
\(v_{s+d}-\sum_{j=1}^d c_jv_{s+d-j}=0\).
The accompanying \texttt{siblings.py} reconstructs each matrix from
the finite-set definitions above and checks these vector equalities.
Thus the finite premises, including every coordinate, are explicit.

Multiplying~\eqref{eq:row-cert} by
\(\alpha_h Q_h^{n-1-s-d}\) proves the row recurrence for every
\(n\ge s+d+1=N\), by~\eqref{eq:generic-row}.
The same argument applied to~\eqref{eq:generic-column} proves the
column recurrences for heights at least 3, 5 and 5, respectively.
For column 6, the additional case \(h=4\) is
\[
 842616=25(40392)-90(1944)+81(96),
\]
so its stated range is \(h\ge4\). In particular, no conclusion
for all heights depends on extrapolating a finite scalar prefix.

For four rows the first seven counts are
\[
81,384,2592,17496,121176,842616,5871528.
\]
Multiplying their power series by \(1-9z+15z^2-6z^3\) gives the coefficients
\[
 (81,-345,351,-558,288,-1080,648)
\]
in degrees 1--7.
The recurrence makes every higher coefficient zero. This proves the
generating function; the nonzero coefficient 648 in degree 7 proves
the sharp starting index.
\end{proof}

\noindent The source coverage is as follows. ``Proved'' refers to an
empirical assertion at the specified FORMULA locator. The larger-strip
proofs use the finite integer identities above. No minimality claim is
made for their recurrence degrees. ``Reproved'' credits the older grid
enumeration together with Theorem~\ref{thm:boundary}.
\begin{center}
\small
\begin{tabularx}{\textwidth}{@{}p{.25\textwidth}XX@{}}
\toprule
Source & Statement & Coverage\\\midrule
A233156, FORMULA 1 & Hardin's recurrence for \(n>4\) & Reproved; prior proof Biggs--Meredith (1976), boundary reduction in Theorem~\ref{thm:boundary}\\
A233156, FORMULA 2 & Barker's generating function & Reproved; prior grid series Ro\v{c}ek--Shrock--Tsai (1997), full boundary numerator in Section~\ref{sec:boundary}\\
A233155, row 3 & The same recurrence & Reproved; prior proof Biggs--Meredith (1976)\\
A233155, rows 1, 2 & Orders one with the stated initial ranges & Proved, Corollary~\ref{cor:thin}\\
A233155, columns 1--4 & Ratios \(3,4,6,9\) as height increases & Proved, Corollary~\ref{cor:thin}\\
A233155, columns 5, 6 & Recurrences of degrees 2 and 3 & Proved, Corollary~\ref{cor:all}; related strip denominators Ro\v{c}ek--Shrock--Tsai (1997)\\
A233155, column 7 & Recurrence of degree 4 & Proved, Corollary~\ref{cor:all}\\
A233155, row 4 & Degree 3 with \(n\ge8\) & Proved, Corollary~\ref{cor:all}; square-strip denominator Ro\v{c}ek--Shrock--Tsai (1997)\\
A233155, rows 5--7 & Degrees 7, 9 and 27 at the stated ranges & Proved, Corollary~\ref{cor:all}\\
A233157, FORMULA 1--2 & Hardin's recurrence and Barker's generating function & Proved, Corollary~\ref{cor:all}\\
A233158--A233160, FORMULA & Explicit higher-row recurrences & Proved, Corollary~\ref{cor:all}\\
General transfer formulas & Equality-pattern transfer for every fixed strip & Not posed; proved, (\ref{eq:generic-row})--(\ref{eq:generic-column})\\
Array-to-grid boundary identity and sequence links & General-color factor and A020698/A052913 identities & Not posed; proved, Theorem~\ref{thm:boundary} and Corollary~\ref{cor:oeis}\\
\bottomrule
\end{tabularx}
\end{center}
The closed form, ratio limit and three-row minimality statement are
not posed in those entries; they are proved consequences. The parent
table's cross-references A000244, A002023, A000400, A055275, A003945
and A005051 identify the elementary sequences in Corollary~\ref{cor:thin};
their other interpretations are outside the present scope. The six-color counting problem
in A233258 has a different alphabet and normalization and is not covered.

\section{Mathematical verification}
The formal counterpart of~\eqref{eq:recurrence} is
\href{https://github.com/the-omega-institute/trureturing/blob/e135d8e2afafffb707a541357df302c35e90baff/D5/S3/Combinatorics/TraceAvoidingArrays/HardinA233156Row.lean}{\texttt{HardinA233156Row.result}}:
for every natural number \(m\),
\[
 U_3(m+5)=5U_3(m+4)-2U_3(m+3),
\]
as an integer identity, with \(U_3(n)=\counts_3(n)\).
The alternating-column bijection is expressed by
\href{https://github.com/the-omega-institute/trureturing/blob/e135d8e2afafffb707a541357df302c35e90baff/D5/S3/Combinatorics/TraceAvoidingArrays/HardinA233155ColumnTransfer.lean}{\texttt{properEquiv}}
and the five-state count for widths at least two by
\href{https://github.com/the-omega-institute/trureturing/blob/e135d8e2afafffb707a541357df302c35e90baff/D5/S3/Combinatorics/TraceAvoidingArrays/HardinA233258PairTransfer.lean}{\texttt{count\_transfer}}.
These declaration names refer to the three-color definitions in the linked
files. The three-row generating function, minimality, closed form and
thin-strip formulas have the written proofs given above.
Corollary~\ref{cor:all} has the written counting reduction and the
specified finite integer identities; no additional formal scope is
attributed to these declarations. In particular, the helper filename
containing \texttt{A233258} refers here to its three-color definitions,
not to the different OEIS entry.

The chromatic-polynomial boundary identity has the written proof of
Theorem~\ref{thm:boundary}. The accompanying \texttt{boundary.py} checks
its extension counts for each proper three-entry boundary column with
two through five colors, compares the three-color array and grid counts,
and checks the cited strip-denominator specializations; the finite
ranges are in \texttt{boundary-verification.json}. This identity and its
general-color version are not posed in the cited OEIS entries.

The accompanying \texttt{verify.py} uses exact integers to check the
column bijection, quotient counts and vector identity, and compares the
27-state literal counts with the five-state counts. Its detailed finite
ranges and results are in \texttt{verification.json}. The all-length
conclusion follows from~\eqref{eq:count}--\eqref{eq:identity}, not from
matching a finite prefix.
The larger-strip matrices and vector identities are covered by
\texttt{siblings.py}; its finite ranges are in
\texttt{siblings-verification.json}. These finite identities, combined
with the written induction, establish Corollary~\ref{cor:all}.

\begin{thebibliography}{12}
\small
\bibitem{source}
R.~H. Hardin, with a generating-function contribution by C.~Barker,
\emph{A233156}, The On-Line Encyclopedia of Integer Sequences,
created 5 December 2013, revision 8, 9 October 2018.
\url{https://oeis.org/A233156}.
The fixed sequence export is
\href{https://raw.githubusercontent.com/oeis/oeisdata/571197984aeda7b8b680fe3a1f6a43f9898cd441/seq/A233/A233156.seq}{\texttt{A233156.seq}}.
\bibitem{parent}
R.~H. Hardin, \emph{A233155}, The On-Line Encyclopedia of Integer
Sequences, created 5 December 2013, revision 6, 3 September 2022.
\url{https://oeis.org/A233155}.
Its fixed export is
\href{https://raw.githubusercontent.com/oeis/oeisdata/571197984aeda7b8b680fe3a1f6a43f9898cd441/seq/A233/A233155.seq}{\texttt{A233155.seq}},
FORMULA entries for columns 1--7 and rows 1--7.
\bibitem{four}
R.~H. Hardin, with a generating-function contribution by C.~Barker,
\emph{A233157}, The On-Line Encyclopedia of Integer Sequences,
revision 8, 9 October 2018, FORMULA lines 1--2.
\href{https://raw.githubusercontent.com/oeis/oeisdata/571197984aeda7b8b680fe3a1f6a43f9898cd441/seq/A233/A233157.seq}{Fixed sequence export}.
\bibitem{higher}
R.~H. Hardin, \emph{A233158, A233159, A233160}, The On-Line
Encyclopedia of Integer Sequences, revision 4, 5 December 2013,
FORMULA line 1 in each entry.
\href{https://raw.githubusercontent.com/oeis/oeisdata/571197984aeda7b8b680fe3a1f6a43f9898cd441/seq/A233/A233158.seq}{A233158};
\href{https://raw.githubusercontent.com/oeis/oeisdata/571197984aeda7b8b680fe3a1f6a43f9898cd441/seq/A233/A233159.seq}{A233159};
\href{https://raw.githubusercontent.com/oeis/oeisdata/571197984aeda7b8b680fe3a1f6a43f9898cd441/seq/A233/A233160.seq}{A233160}.
\bibitem{ac}
P.~Flajolet and R.~Sedgewick, \emph{Analytic Combinatorics},
Cambridge University Press, 2009, Section V.6, ``Transfer matrix models,''
and Figure V.20, p.~356. \url{https://doi.org/10.1017/CBO9780511801655}.
\bibitem{bm}
N.~L. Biggs and G.~H.~J. Meredith, \emph{Approximations for chromatic
polynomials}, Journal of Combinatorial Theory, Series B \textbf{20}
(1976), 5--19. Section~3, Theorem~2, p.~10, and the double-strip
matrix, p.~11. \url{https://doi.org/10.1016/0095-8956(76)90062-9}.
An author-hosted copy is listed at
\url{https://normanbiggs.com/research-papers-to-1988}.
\bibitem{rst}
M.~Ro\v{c}ek, R.~Shrock and S.-H. Tsai, \emph{Chromatic polynomials
for families of strip graphs and their asymptotic limits}, Physica A
\textbf{252} (1998), 505--546.
\url{https://doi.org/10.1016/S0378-4371(98)00034-X}.
Preprint, 13 December 1997:
\url{https://arxiv.org/abs/cond-mat/9712148v1}.
\bibitem{ss}
J.~Salas and A.~D. Sokal, \emph{Transfer matrices and partition-function
zeros for antiferromagnetic Potts models I. General theory and
square-lattice chromatic polynomial}, Journal of Statistical Physics
\textbf{104} (2001), 609--699.
\url{https://doi.org/10.1023/A:1010376605067}.
Preprint, first posted 19 April 2000, version 2 dated 8 January 2001:
\url{https://arxiv.org/abs/cond-mat/0004330v2}.
Section~5.2, eqs.~(5.3)--(5.4) and footnote~28.
\bibitem{gh}
J.~Ginepro and T.~C. Hull, \emph{Counting Miura-ori foldings}, Journal
of Integer Sequences \textbf{17} (2014), Article 14.10.8, published
5 November 2014. Section~3, p.~6, and Theorem~4, p.~11.
\url{https://cs.uwaterloo.ca/journals/JIS/VOL17/Hull/hull.html}.
\bibitem{spahn}
G.~Spahn, \emph{Counting Classes of Matrices and More Using Experimental
Mathematics}, Ph.D. dissertation, Rutgers, The State University of New
Jersey, 4 March 2024. Section~6.1, pp.~44--45.
\url{https://sites.math.rutgers.edu/~zeilberg/Theses/GeorgeSpahnThesis.pdf}.
\bibitem{related-one}
D.~W. Wilson and contributors, \emph{A020698}, The On-Line Encyclopedia
of Integer Sequences, revision 62, 8 September 2022, NAME and FORMULA;
grid-coloring comment by A.~Woods, 7 September 2011.
\href{https://raw.githubusercontent.com/oeis/oeisdata/571197984aeda7b8b680fe3a1f6a43f9898cd441/seq/A020/A020698.seq}{Fixed sequence export}.
\bibitem{related-two}
The On-Line Encyclopedia of Integer Sequences, \emph{A052913},
revision 81, 8 September 2022, NAME and FORMULA.
\href{https://raw.githubusercontent.com/oeis/oeisdata/571197984aeda7b8b680fe3a1f6a43f9898cd441/seq/A052/A052913.seq}{Fixed sequence export}.
\end{thebibliography}

\bigskip
\noindent\small \copyright\ 2026 Haobo Ma and Wenlin Zhang.
Licensed under Creative Commons Attribution 4.0 International
(CC BY 4.0): \url{https://creativecommons.org/licenses/by/4.0/}.
\end{document}
